
%%%%%%%%%%%%%%%%%%%%%%%%%%%%%%%%%%%%%%%%
%           main body
%%%%%%%%%%%%%%%%%%%%%%%%%%%%%%%%%%%%%%%%


%-----------------------------------
%	TITLE PAGE
%-----------------------------------

\title[Optimization Basics]{\LARGE \bfseries \quad 优化基础知识} % The short title appears at the bottom of every slide, the full title is only on the title page
% \author[Basic Concepts] % Your institution as it will appear on the bottom of every slide, may be shorthand to save space
% \author{Singular Value Decomposition} % Your name
% \date{\today} % Date, can be changed to a custom date
\date{}

%------------------------------------------------

\begin{document}


%------------------------------------------------

\begin{frame}

\titlepage % Print the title page as the first slide

\end{frame}

%------------------------------------------------

% \begin{frame}
% \frametitle{内容提要} % Table of contents slide, comment this block out to remove it
% \tableofcontents % Throughout your presentation, if you choose to use \section{} and \subsection{} commands, these will automatically be printed on this slide as an overview of your presentation
% \end{frame}

%-----------------------------------
%	PRESENTATION SLIDES
%-----------------------------------

\section{基本概念}

%------------------------------------------------

\begin{frame}{数学优化问题，{\smaller Mathematical Optimization Problem}}

% \begin{block}
数学优化问题，或简称优化问题，指的是求以下的量的问题
$$\min\limits_{\mathbf{x}} \left\{ f_0(\mathbf{x}) \ \middle|\ f_i(\mathbf{x}) \le (\text{或} = ) b_i, \; i = 1,\ldots,m \right\}$$
其中
\begin{itemize}
\item $\mathbf{x} = (x_1,\ldots,x_n)^T$ 为优化变量（Optimization Variable）
\item $f_0: \mathbb{R}^n \to \mathbb{R}$ 被称作目标函数（Objective Function），其定义域我们一般记作$\mathcal{D}$或$\operatorname{dom}f$
\item $f_i: \mathbb{R}^n \to \mathbb{R}, \; i = 1,\ldots,m,$ 被称作约束函数（Constraint Functions）。如果一个优化问题没有约束函数，则被称作一个无约束问题。
\item $b_1,\ldots,b_m$ 被称作约束边界（Bounds）
\end{itemize}
% \end{block}

\end{frame}

%------------------------------------------------

\begin{frame}{最优解，Optimal Solution}

在一个优化问题中，满足所有约束条件
$$f_i(\mathbf{x}) \le b_i,\; i = 1,\ldots,m,$$
的向量$\mathbf{x}$被称作可行的（Feasible）。全体可行点组成的集合被称作可行集。

\vspace{1em}

如果可行集中一个向量$\mathbf{x}^*$, 在目标函数$f_0$上取值最小，即
$$\mathbf{x}^* = \argmin\limits_{\mathbf{x} \in \mathcal{D}} \left\{ f_0(\mathbf{x}) \ \middle|\ f_i(\mathbf{x}) \le b_i, \; i = 1,\ldots,m \right\}$$
则称$\mathbf{x}^*$为这个优化问题的最优解。相应地，目标函数$f_0$在最优解$\mathbf{x}^*$处的取值$p^* = f_0(\mathbf{x}^*)$被称作最优值。
% \end{block}

\end{frame}

%------------------------------------------------

\begin{frame}{一些基本的优化问题}

\begin{itemize}
\item 最小二乘法：
$$\min\limits_{\mathbf{x}} \Vert A\mathbf{x} - \mathbf{b} \Vert_2^2 = \sum (\mathbf{a}_i^T\mathbf{x} - b_i)^2$$
这是一个无约束问题。
\vspace{1em}
\item 线性规划 (Linear Programming)
$$\min\limits_{\mathbf{x}} \left\{ \mathbf{c}^T\mathbf{x} \ \middle|\ \mathbf{a}_i^T\mathbf{x} \le b_i, \; i = 1,\ldots,m \right\}$$
\end{itemize}

\end{frame}

%------------------------------------------------

\section{基础算法}

%------------------------------------------------

\begin{frame}{无约束优化}

% \begin{block}
无约束优化问题具有以下形式
$$\text{minimize} \quad f(\mathbf{x})$$
如果$f$是凸函数 (convex function)，那么
$$\mathbf{x}^* \text{ 是最优解 } \Longleftrightarrow \nabla f(\mathbf{x}^*) = \mathbf{0}$$
这个方程（组）一般没有解析解，必须采用迭代算法求解，即寻找所谓的极小化点列
$$\mathbf{x}^{(0)}, \ldots, \mathbf{x}^{(k)}, \ldots,$$
使得当$k\to\infty$时，$f(\mathbf{x}^{(k)}) \to p^*$
% \end{block}

\end{frame}

%------------------------------------------------

\begin{frame}{下降方法}

极小化点列将按以下方式产生
$$\mathbf{x}^{(k+1)} = \mathbf{x}^{(k)} + t^{(k)} \Delta \mathbf{x}^{(k)}, \text{ 且有 } f(\mathbf{x}^{(k+1)}) < f(\mathbf{x}^{(k)})$$
其中
\begin{itemize}
\item $\Delta \mathbf{x}^{(k)}$被称为第$k$次迭代搜索方向，或步径
\item $t^{(k)} > 0$被称作第$k$次迭代步长 (step size)，又称学习率 (learning rate)
\end{itemize}

\end{frame}

%------------------------------------------------

\begin{frame}{下降方法}

\begin{algorithm}[H]
\SetAlgoLined
\DontPrintSemicolon
\SetKwInOut{Input}{Input}
\SetKwInOut{Output}{Output}
\SetKwRepeat{Do}{do}{while}

给定初始点$\mathbf{x}$ \;
\Do{达到停止条件}
{
确定下降方向$\Delta \mathbf{x}$ \;
搜索并选择步长$t$ \;
$ \mathbf{x} \gets \mathbf{x} + t \Delta \mathbf{x}$ \;
}
\Output{近似最优解$\mathbf{x}$}
\caption{通用下降方法}
\end{algorithm}

\vspace{1em}
\pause

\begin{itemize}
\item $\Delta \mathbf{x} = - \nabla f(\mathbf{x}) \; \rightsquigarrow \; \text{\color{red} 梯度下降法}$
\item $\Delta \mathbf{x}_{nsd} = \argmin\limits_{\mathbf{v}} \{ \nabla f(\mathbf{x})^T \mathbf{v} \ |\ \Vert \mathbf{v} \Vert = 1 \} \; \rightsquigarrow \; \text{\color{red} 最速下降法}$
\item $\Delta \mathbf{x}_{nt} = - H(f)(\mathbf{x})^{-1} \nabla f(\mathbf{x}) \; \rightsquigarrow \; \text{\color{red} Newton法}$
\end{itemize}

\end{frame}

%------------------------------------------------
% % closing page
% \begin{frame}
% \HUGE{\centerline{\bfseries {\color{gray} The} {\color{zkblue} End}}}

% \pause

% \vspace{1.2em}

% \Large{\centerline{\bfseries {\color{gray}Thank} \quad {\color{zkblue} You}}}

% \end{frame}

%----------------------------------------------------------------------------------------

\end{document}
